\chapter{Conditionnement, indépendance et Bayes}
\section*{Exercice 19.1 -- Trois probabilités}
$P(A)=0{,}4$, $P(B)=0{,}5$, $P(A\cap B)=0{,}3$. Calculer $P_A(B)$, $P_B(A)$ et $P(A\cup B)$.

\section*{Exercice 19.2 -- Probabilités totales}
Une machine A produit $60\%$ des pièces avec $2\%$ de défauts ; B produit le reste avec $5\%$ de défauts. Calculer la probabilité d'un défaut.

\section*{Exercice 19.3 -- Bayes}
Dans l'exercice précédent, sachant qu'une pièce est défectueuse, quelle est la probabilité qu'elle provienne de B ?

\section*{Exercice 19.4 -- Tester l'indépendance}
Deux événements ont probabilités $1/3$ et $3/5$, leur intersection $1/5$. Sont-ils indépendants ? Qu'en est-il de leurs complémentaires ?

\section*{Exercice 19.5 -- Incompatibilité}
$A$ et $B$ sont incompatibles et ont des probabilités strictement positives. Prouver qu'ils sont dépendants.

\section*{Exercice 19.6 -- Événement nul}
Si $P(A)=0$, peut-on définir $P_A(B)$ par la formule habituelle ? $A$ est-il indépendant de $B$ ?
